% Standalone source-record documentation; no research attempts.
\documentclass[11pt]{article}
\usepackage[margin=25mm]{geometry}
\usepackage{fontspec,unicode-math}
\setmainfont{Latin Modern Roman}
\setmathfont{Latin Modern Math}
\usepackage{amsmath,amssymb,mathtools,xurl,hyperref}
\hypersetup{colorlinks=true,urlcolor=blue,linkcolor=blue}
\setcounter{secnumdepth}{0}
\setlength{\parindent}{0pt}
\setlength{\parskip}{5pt}
\setlength{\emergencystretch}{3em}
\newcommand{\sref}[2]{\hyperlink{source:#1:#2}{[S#2]}}
\newcommand{\cataloguescope}[1]{\paragraph{Source record.}#1}
\newcommand{\floor}[1]{\left\lfloor#1\right\rfloor}
\begin{document}
\section{Better-Than-Two Approximation for Multiplicative Connectivity Augmentation}\label{statement}

\subsection{Definitions and mathematical statement}

\paragraph{Shared notation.}Standard number systems are $\mathbb N,\mathbb Z,\mathbb Q,\mathbb R,\mathbb C$. Specialized notation and index ranges follow the cited source and the definitions below.


In $\alpha$-multiplicative connectivity augmentation the input is a $k$-edge-connected graph and a set of candidate links with nonnegative costs. One selects minimum-cost links so that the augmented graph is at least $\lceil\alpha k\rceil$-edge-connected. Edge connectivity is the minimum number of edges in a nontrivial cut. An approximation factor below two means a polynomial-time algorithm with cost at most $\rho$ times optimum for some $\rho<2$.
\paragraph{Record scope and qualifications.}
The catalogue title is retained as metadata. Named certificates, contracts, bounds or frameworks in that title are not asserted as results of the original source extract.
The supplied registry flags this extraction as unrecoverable. It is displayed as a literal source fragment; the definitions provide context rather than a replacement formal statement.

Source-record limitation: Source flags the statement as unrecoverable: The official report confirms that this is a duplicate extraction of the multiplicative connectivity augmentation problem retained in OWR-8415356-018. No separate clean mathematical statement is supplied; this source ID preserves its provenance.

Related catalogue records: \href{https://mehmetmars7.github.io/Erdosproblems-llm-hunter/problem.html?type=open_problems&id=30004907}{30004907}. Separate source IDs are preserved for their provenance.
The source text below is rendered literally where necessary; damaged formulas and composite extracts are not silently repaired into a new mathematical claim.
The supplied catalogue marks this record as unrecoverable. Its extract and terminology are documented, but a complete original problem has not been reconstructed.
\paragraph{Statement recorded in the supplied source.}
Therefore, we propose first studying \textbackslash{}alpha-multiplicative connectivity augmentation for \textbackslash{}alpha \textbackslash{}leq 6/5 which may alternately be called polygon augmentation, following the naming convention of the problems tree augmentation and cactus augmentation in which the goal is to increase the connectivity of a graph by 1. The main open problem is to design a better-than-2 approximation for this problem for any constant \textbackslash{}alpha > 1. Open Problem: \textbackslash{}alpha-Multiplicative Connectivity Augmentation and Polygon Augmentation Nathan Klein \sref{statement}{2}
\paragraph{Source status.}
The supplied catalogue labels this source record partially solved; that metadata was checked against the supplied snapshot on 2026-10-11. The label is preserved for provenance and is not an independent verification of present open status or mathematical correctness. The catalogue literature assessment is dated 2026-08-22.
\subsection{Short English statement}
The supplied fragment concerns better-Than-Two Approximation for Multiplicative Connectivity Augmentation. Its complete mathematical target is not recovered.
\subsection{Sources}
\begin{itemize}
\item[\hypertarget{source:statement:1}{S1}] Ulam AI and the UnsolvedMath Contributors, UnsolvedMath: A Curated Collection of Open Mathematics Problems, catalogue entry. Catalogue curation and metadata: CC BY 4.0. The source extract is retained; definitions, English summaries and qualifications are editorial additions.
\url{https://huggingface.co/datasets/ulamai/UnsolvedMath}.
\item[\hypertarget{source:statement:2}{S2}] Combinatorial Optimization (2022), Oberwolfach Reports, DOI 10.4171/owr/2021/53, pp. 2893--2954
\url{https://doi.org/10.4171/owr/2021/53}.
\item[\hypertarget{source:statement:3}{S3}] Weighted connectivity augmentation
\url{https://arxiv.org/abs/2209.07860}.
\item Source-specific Oberwolfach Reports licences, including CC BY-SA 4.0 where applicable, are retained; see THIRD\_PARTY\_NOTICES.md.
\end{itemize}
\paragraph{Preparation and review.}
The original source and UnsolvedMath contributors retain attribution for the source extract. Codex (GPT-6; exact serving revision not exposed) assisted with definitions, English summaries, scope notes and formatting. No new human mathematical authorship or independently verified solution is asserted. Catalogue identity checked against the complete supplied registry and current repository. Distinct upstream records, including repeated questions and broad or damaged extracts, are retained for complete coverage of the source CS category. A repeated record is not a new independent conjecture. Complete mathematical-equivalence and current-literature review has not been performed. Related canonical identities: 30004907. Preparation and final review follow contribution prompts C and D, with source gaps disclosed. Structural and rendering checks do not establish a new theorem or open status.
\label{end:statement}
\end{document}
